\documentclass[11pt]{article}
\usepackage[margin=0.8in]{geometry}
\usepackage[T1]{fontenc}
\usepackage{lmodern,amsmath,amssymb,amsthm,booktabs,array}
\usepackage[colorlinks=true,linkcolor=blue,urlcolor=blue,citecolor=blue]{hyperref}
\hypersetup{pdftitle={Odd-degree irreducibility, a word identity, and arithmetic assertions},pdfauthor={Blake Stone},pdfsubject={Complete proofs, source correspondence, and Lean verification}}
\setlength{\parindent}{0pt}
\setlength{\parskip}{0.4em}
\newtheorem{theorem}{Theorem}[section]
\newtheorem{lemma}[theorem]{Lemma}
\newtheorem{corollary}[theorem]{Corollary}
\newcommand{\N}{\mathbb N}
\newcommand{\Z}{\mathbb Z}
\newcommand{\Q}{\mathbb Q}
\newcommand{\R}{\mathbb R}
\begin{document}
\begin{center}
{\LARGE\bfseries Odd-degree irreducibility, a word identity,\\[3pt]
and arithmetic assertions}\par
\vspace{0.5em}
{\large Blake Stone}\par
Proof developed with OpenAI Astra 6.2\par
November 27, 2026
\end{center}

\textbf{Abstract.}
We prove that every even-order derivative of
$(1-x)/(1-x-x^2)$ has irreducible reduced numerator over $\Q$,
settling the odd-degree clause recorded in A326925. The normalized
polynomials are in fact irreducible over $\Q(\sqrt5)$. The proof reduces
them by a linear-fractional transformation to odd binomials and uses an
integer-trace obstruction in the golden-ratio field. We also prove an exact
concatenation identity for the A288665 replacement process, establishing
its conjectured A288668 length recurrence, and refute the A134559
antidiagonal-product characterization with the prime 54217. Finally, a
uniform polynomial congruence proves the three prime-shift assertions
recorded in A000033, A000159, and A000181. Relevant older congruence
literature could not be fully inspected, so historical novelty is not
established for those three assertions. Four independently audited Lean
packages contain complete source-definition bridges and proof terms for
the stated results: five proved assertions and one refutation. No claim
of exhaustive historical priority is made.

\begin{center}
\begin{tabular}{@{}l>{\raggedright\arraybackslash}p{4.8in}@{}}
\toprule
Source & Precisely established contribution\\
\midrule
A326925 & Odd-degree irreducibility; the even-degree factorization clause remains outside this result.\\
A288665/A288668 & One shared length assertion, proved through a stronger identity of actual words.\\
A134559 & One characterization refuted by a fully checked finite counterexample.\\
A000033/159/181 & Three prime-shift assertions proved and formalized; historical novelty unverified.\\
\bottomrule
\end{tabular}
\end{center}

\section{A326925: every odd-degree numerator is irreducible}

Clark Kimberling's A326925~\cite{A326925} records polynomials beginning
\[
\begin{aligned}
P_1&=1-x,& P_2&=2x-x^2,\\
P_3&=1+3x^2-x^3,& P_4&=1+4x+4x^3-x^4.
\end{aligned}
\]
The entry conjectures that its odd-degree polynomials are irreducible.
Its title shifts the derivative index by one relative to the degree;
the displayed polynomials and executable program fix the intended
convention. We index by degree and prove the original derivative
correspondence explicitly. The program's nonzero sign and factorial
normalizations are irrelevant to irreducibility over a field.

\begin{theorem}\label{thm:odd}
Let $f(x)=(1-x)/(1-x-x^2)$. For every even integer $n\geq0$, the reduced
numerator of $f^{(n)}$ is irreducible over $\Q$, of degree $n+1$.
Equivalently, $P_N$ is irreducible for every positive odd degree $N$.
Moreover, $P_N$ remains irreducible over $\Q(\sqrt5)$.
\end{theorem}

The even-degree clause, which predicts two irreducible factors of equal
degree, is a separate assertion and is not proved here.

\subsection{Recurrence and exact derivative correspondence}

Put $D=1-x-x^2$, $A=2x+1$, $B=x^2+x-1=-D$, and introduce $P_0=-1$.
The polynomials satisfy
\begin{equation}\label{eq:recurrence8}
P_{N+2}=AP_{N+1}-BP_N,\qquad P_0=-1,\quad P_1=1-x.
\end{equation}
This recurrence can be used as an integral polynomial definition. It
implies the differential identity
\begin{equation}\label{eq:diff8}
D P_N'+N A P_N=N P_{N+1}.
\end{equation}
For a direct induction, use $A'=2$, $B'=A$, and $D=-B$. Substituting
the recurrence at indices $N$ and $N+1$ in the desired identity at
$N+2$ reduces its difference from zero to $A$ times the identity at
$N+1$ minus $B$ times the identity at $N$. The base indices 0 and 1
follow by substitution.

The quotient rule and~\eqref{eq:diff8} now give, by induction on $n$,
\begin{equation}\label{eq:derivative8}
\boxed{f^{(n)}(x)=\frac{n!P_{n+1}(x)}{D(x)^{n+1}}.}
\end{equation}
Indeed, differentiating the right side gives numerator
$n![D P_{n+1}'+(n+1)A P_{n+1}]=(n+1)!P_{n+2}$ over $D^{n+2}$.
The formal development defines rational differentiation using canonical
numerator and denominator, and proves the quotient rule for every
representation $p/q$ with $q\ne0$. Thus the bridge is an equality of
rational functions, not a recurrence assumed for a conjectured numerator.

\subsection{Closed form in a quadratic field}

Set
\[
u=\frac{1+\sqrt5}{2},\qquad v=1-u,\qquad K=\Q(u).
\]
Then $u+v=1$, $uv=-1$, $u^2=u+1$, and $1<u<2$. The characteristic
roots of~\eqref{eq:recurrence8}, as a recurrence in $N$, are $x+u$
and $x+v$. The initial values therefore give
\begin{equation}\label{eq:closed8}
(u-v)P_N(x)=v^2(x+u)^N-u^2(x+v)^N.
\end{equation}
This identity also follows by checking $N=0,1$ and the recurrence
directly, without adjoining any further elements. Since $v^2u^4=u^2$,
\begin{equation}\label{eq:mobius8}
(u-v)P_N(x)=v^2\big((x+u)^N-u^4(x+v)^N\big).
\end{equation}
The leading coefficient on the right is $v^2(1-u^4)\ne0$, so
$\deg P_N=N$.

\begin{lemma}\label{lem:unit8}
For every odd prime $p$, the element $u^4$ is not a $p$-th power in $K$.
\end{lemma}
\begin{proof}
Suppose $b^p=u^4$ with $b\in K$. The element $u$ is integral because it
satisfies the monic equation $u^2-u-1=0$. The integrality of $b^p$ implies
that of $b$. The field norm gives
\[
\operatorname{N}_{K/\Q}(u)=-1,\qquad
\operatorname{N}_{K/\Q}(b)^p=1.
\]
Since $p$ is odd and this norm is rational, $\operatorname{N}(b)=1$.
In the real embedding under consideration, $b^p=u^4>1$, hence $b>1$.
Also $u^4=3u+2<8$. Since $p\geq3$, the possibility $b\geq2$ would
give $b^p\geq8$, a contradiction. Thus $1<b<2$.

The number $b$ is irrational, because $b^p=u^4=3u+2$ is irrational.
Its other conjugate is $1/b$, by the norm identity. The trace of an
algebraic integer in this quadratic field is an integer, but
\[
2<b+\frac1b<3\qquad(1<b<2).
\]
This contradiction proves the lemma.
\end{proof}

The Lean proof derives all the needed degree, norm, trace, and
integrality facts from the explicit quadratic polynomial. In particular,
it does not assume a fundamental-unit classification of $K$.

\subsection{Irreducibility and reduction}

The classical odd binomial irreducibility criterion states that
$Y^N-c$ is irreducible over a characteristic-zero field when $N$ is odd
and $c$ is not a $p$-th power for any prime $p\mid N$. This criterion
is proved in the pinned mathlib dependency~\cite{mathlib}. By
Lemma~\ref{lem:unit8},
\begin{equation}\label{eq:binomial8}
Y^N-u^4\text{ is irreducible over }K\quad(N\text{ odd}).
\end{equation}

\begin{lemma}\label{lem:mobius8}
Let $F$ be a field, let $a,b,c\in F$ satisfy $a\ne b$ and $c\ne1$,
and suppose $Y^N-c$ is irreducible. Then
\[
Q(X)=(X-a)^N-c(X-b)^N
\]
is irreducible over $F$.
\end{lemma}
\begin{proof}
The hypothesis forces $N>0$, and $c\ne1$ ensures $\deg Q=N$.
Choose an irreducible factor $g$ of $Q$ and let $x$ be its root in the
field $E=F[X]/(g)$. The equation $Q(x)=0$ shows $x\ne b$, since
otherwise $(b-a)^N=0$. Hence $y=(x-a)/(x-b)\in E$ satisfies $y^N=c$.
Its minimal polynomial is the monic irreducible polynomial $Y^N-c$,
so
\[
N\leq[E:F]=\deg g\leq\deg Q=N.
\]
Thus $g$ has the full degree of $Q$ and is associated to it. This proves
irreducibility.
\end{proof}

Apply Lemma~\ref{lem:mobius8} with $a=-u$, $b=-v$, and $c=u^4$.
Equations~\eqref{eq:mobius8} and~\eqref{eq:binomial8} give
irreducibility of $P_N$ over $K$ for odd $N$. It follows over $\Q$
as well, since a factorization over $\Q$ would persist over $K$.

For odd $N\geq3$, the irreducible polynomial $P_N$ cannot divide $D$
because $\deg P_N=N>2$. For $N=1$, $P_1=1-x$ does not divide $D$
because $D(1)=-1$. Therefore no factor of $P_N$ cancels against a power
of $D$ for odd $N$. In~\eqref{eq:derivative8}, the factor $n!$ is a
unit of $\Q$. For even $n$, the reduced numerator is consequently
associated to $P_{n+1}$, completing the proof of Theorem~\ref{thm:odd}.

\section{A288665/A288668: a literal word identity}

Start with $w_0=00$. Obtain $w_{n+1}$ by scanning $w_n$ from left to
right and making the nonoverlapping replacements
\[
00\longmapsto0110,\qquad10\longmapsto000,
\]
copying any letter that does not start a matching pair. These are the
literal source rules in A288665~\cite{A288665}. A288668~\cite{A288668}
conjectures the order-three recurrence for the resulting lengths.

\begin{theorem}\label{thm:word8}
For every $n\geq0$,
\[
\boxed{w_{n+3}=w_nw_nw_{n+1},}
\]
where juxtaposition denotes concatenation. In particular, if $L_n=|w_n|$,
\[
L_{n+3}=L_{n+1}+2L_n,\qquad (L_0,L_1,L_2)=(2,4,5).
\]
\end{theorem}
\begin{proof}
Introduce three token codes
\[
A=00,\qquad B=0110,\qquad C=01000.
\]
One pass gives $A\mapsto B$, $B\mapsto C$, and $C\mapsto AAB$.
For example,
\[
0110\mapsto0\,1\,000=01000,
\qquad01000\mapsto0\,000\,0110=00\,00\,0110.
\]
Each token's final pair is consumed by its rewrite, so the same identity
holds when any binary suffix follows the token. Thus there is no
boundary-crossing match between these token codes. This arbitrary-suffix
property is proved explicitly in Lean.

Let $\phi$ be the substitution on token words defined by
$\phi(A)=B$, $\phi(B)=C$, $\phi(C)=AAB$. Decoding $\phi^n(A)$ therefore
gives the exact literal iterate $w_n$, by induction. Since $\phi^3(A)=AAB$,
applying $\phi^n$ and decoding yields the stated concatenation identity.
The first words are $00$, $0110$, and $01000$, so their lengths are
$2,4,5$. Taking lengths completes the proof.
\end{proof}

The source prose gives inconsistent labels for the three initial
indices. Its data, offset, generating function, and program all agree
with the values at indices $0,1,2$ used above. This clerical error is
documented, not counted as another mathematical refutation. The shared
length assertion is counted once even though both entries record it.

\begin{corollary}
For any assignment of nonnegative integer weights to the two letters,
the total weight $W_n$ satisfies $W_{n+3}=W_{n+1}+2W_n$. Moreover,
for all $n,t\geq0$,
\[
w_n\text{ is a prefix of }w_{n+3t},\qquad
|w_{n+3t}|\geq2^t|w_n|.
\]
\end{corollary}
\begin{proof}
Weights add under concatenation, proving the first claim. The word
identity makes $w_n$ a prefix of $w_{n+3}$ and gives
$|w_{n+3}|\geq2|w_n|$. Iterate these two facts.
\end{proof}

All lengths are positive. The nested words in each residue class modulo
three therefore have unbounded lengths and define the three limiting
words discussed in the related source entries. Their formal prefix and
growth results are included as consequences, not separately counted
conjecture resolutions.

\section{A134559: an explicit counterexample}

Lechoslaw Ratajczak's 2019 comment in A134559~\cite{A134559} defines
\[
b_n(k)=\sum_{\substack{1\leq u\leq n\\n\bmod u=k}}u,
\qquad
P(n)=\prod_{k=0}^{\lceil n/2\rceil-1}b_n(k).
\]
It conjectures that the positive integers satisfying
\begin{equation}\label{eq:gcd8}
\gcd(P(n),n)=1,\qquad \gcd(P(n+1),n+1)=2
\end{equation}
are exactly the primes $p>5$ for which $(p+1)/2$ is prime. This prime
family is A005383 with its first two entries excluded~\cite{A005383}.

\begin{theorem}\label{thm:counter8}
The characterization~\eqref{eq:gcd8} is false. Specifically,
\[
p=54217,\qquad \frac{p+1}{2}=27109
\]
are both prime, but
\[
b_{54217}(2737)=162651=3\cdot54217,
\qquad \gcd(P(54217),54217)=54217.
\]
\end{theorem}
\begin{proof}
Primality of 54217 and 27109 is verified by trial division, or by the
proof-producing primality tactic in the accompanying Lean development.
If $54217\bmod u=2737$, then $u>2737$ and
\[
54217=qu+2737,\qquad qu=51480.
\]
Conversely, every divisor $u$ of 51480 greater than 2737 gives this
remainder. Such a $u$ is at most 51480, so it lies in the original
summation interval. Since $u\geq2738$, its positive complementary divisor
$q$ is at most 18. Enumerating these possible divisors gives exactly
\[
\begin{gathered}
2860,3432,3960,4290,4680,5148,5720,6435,\\
8580,10296,12870,17160,25740,51480.
\end{gathered}
\]
Their sum is $162651=3\cdot54217$. The index 2737 lies in the original
product interval because $2737<\lceil54217/2\rceil=27109$. Thus
$54217\mid P(54217)$ and the first gcd in~\eqref{eq:gcd8} equals
54217, not 1. The second gcd condition cannot remedy this failure.
\end{proof}

The formal proof establishes both inclusions in the displayed remainder
fiber, both primality claims, the exact sum, its membership as a product
factor, and the final negation of the proposed characterization over
positive integers. It never evaluates the enormous full product.
No minimal-counterexample claim or classification of all failures is
made.

\section{Prime-shift congruences for three m\'enage sequences}
\label{sec:menage}

The entries A000033, A000159, and A000181 record three congruence conjectures
of Mark van Hoeij, dated June 10, 2019 \cite{A000033,A000159,A000181}.
We prove all three from a single polynomial identity and formalize the exact
source-definition bridge in Lean. Historical novelty is not established:
Carlitz's and Stevens's earlier papers concern m\'enage congruences, but their
complete theorem statements were unavailable for this review
\cite{Carlitz1952,Carlitz1954,Stevens1981}.

Following A058087, put $U_0(t)=1$ and, for $n\geq1$, define the classical
m\'enage hit polynomial
\begin{equation}\label{eq:menage-source}
 U_n(t)=\sum_{k=0}^{n}\frac{2n}{2n-k}
       \binom{2n-k}{k}(n-k)!(t-1)^k.
\end{equation}
Write $h_d(n)=[t^d]U_n(t)$. The triangle in A058087 lists coefficients in
descending degree, so $h_d(n)=\mathrm{A058087}(n,n-d)$ when $d\leq n$
\cite{A058087}. The relevant fixed degrees are $2,3,4$.

\begin{theorem}\label{thm:menage-shift}
For every prime $p$ and every $n\geq1$,
\begin{equation}\label{eq:menage-shift}
 U_{n+2p}(t)-2(t^p-1)U_{n+p}(t)+(t^p-1)^2U_n(t)
 \equiv0\pmod p
\end{equation}
coefficientwise. Consequently, for every $d<p$,
\[
 p\mid h_d(n)+2h_d(n+p)+h_d(n+2p).
\]
For $p=2$, this divisibility also holds for every $d<4$.
\end{theorem}

\begin{proof}
For $n\geq1$ and $j\geq0$, set
\[
 R(n,j)=\binom{n+j}{2j}+\binom{n+j-1}{2j},
\]
with a binomial coefficient equal to zero when its lower argument exceeds its
upper argument. If $j\leq n$, the identity
\[
 (n+j)\binom{n+j-1}{2j}=(n-j)\binom{n+j}{2j}
\]
shows that $R(n,j)=2n\binom{n+j}{2j}/(n+j)$. Reindexing
\eqref{eq:menage-source} by $j=n-k$ and using binomial symmetry gives the
integral polynomial identity
\begin{equation}\label{eq:menage-integral}
 U_n(t)=\sum_{j=0}^{n}j!R(n,j)(t-1)^{n-j}.
\end{equation}
In particular, all coefficients in the source polynomial are integers.

Work now in $\mathbb F_p[x]$. Frobenius implies
\[
\begin{aligned}
 &(1+x)^{a+2p}-2(1+x)^{a+p}+(1+x)^a\\
 &\qquad=(1+x)^a\bigl((1+x)^p-1\bigr)^2
 =x^{2p}(1+x)^a
\end{aligned}
\]
for every $a\geq0$. Therefore, for $r<2p$,
\begin{equation}\label{eq:menage-binomial}
 \binom{a+2p}{r}-2\binom{a+p}{r}+\binom ar\equiv0\pmod p.
\end{equation}
For $j<p$, apply \eqref{eq:menage-binomial} with $r=2j$, first with $a=n+j$
and then with $a=n+j-1$. Both upper arguments are nonnegative. Adding gives
\begin{equation}\label{eq:menage-rook}
 R(n+2p,j)-2R(n+p,j)+R(n,j)\equiv0\pmod p.
\end{equation}

Terms of \eqref{eq:menage-integral} with $j\geq p$ vanish modulo $p$ because
$p\mid j!$. Also $R(n,j)=0$ for $j>n$. Thus the sums for $n,n+p,n+2p$ may
all be compared over $0\leq j<p$, even when $n<p$. Multiply
\eqref{eq:menage-rook} by $j!(t-1)^{n+2p-j}$ and sum over that range.
The common exponent is nonnegative, and the zero terms with $j>n$ justify
completing the shorter sum. We obtain
\[
 U_{n+2p}(t)-2(t-1)^pU_{n+p}(t)+(t-1)^{2p}U_n(t)=0
 \quad\text{in }\mathbb F_p[t].
\]
Since $(t-1)^p=t^p-1$ in this ring, this is \eqref{eq:menage-shift}.
For $d<p$, taking coefficients of $t^d$ removes every term containing $t^p$
and yields the stated divisibility.

When $p=2$, the middle term is zero and $(t-1)^4=t^4+1$.
Hence $U_{n+4}+(t^4+1)U_n=0$ over $\mathbb F_2$, which gives the same
conclusion in every degree $d<4$.
\end{proof}

\begin{corollary}
Each of the following three source assertions holds on its entire domain:
\[
\begin{array}{c|c|c|c}
\text{Entry}&a(n)&\text{Index domain}&\text{Allowed primes}\\\hline
\mathrm{A000033}&h_2(n)&n\geq1&\text{all primes }p\\
\mathrm{A000159}&h_3(n)&n\geq3&p\ne3\\
\mathrm{A000181}&h_4(n)&n\geq4&p>3
\end{array}
\]
In each case, $p\mid a(n)+2a(n+p)+a(n+2p)$.
\end{corollary}

\begin{proof}
For $d=2$, every odd prime has $d<p$; for $d=3$, every prime other than
$2,3$ has $d<p$; for $d=4$, every prime greater than $3$ has $d<p$.
The remaining prime $2$ in the first two assertions is covered by the
additional $d<4$ conclusion. At $n=1$, A000033's explicit degree-two
coefficient formula gives zero, as required by its source offset.
\end{proof}

The Lean package proves equality between the literal rational rook sum
\eqref{eq:menage-source} and the integral definition
\eqref{eq:menage-integral}, including $U_0=1$. Its final declarations
\texttt{A000033\_prime\_shift}, \texttt{A000159\_prime\_shift}, and
\texttt{A000181\_prime\_shift} assert integer divisibility with exactly the
domains above. No conjectural recurrence or finite verification hypothesis is
used. These are three recorded assertions proved and formalized; the available
source review does not establish that they were historically open.


\section{Formal verification, attribution, and scope}

All four packages pin Lean 4.19.0 and mathlib revision
\texttt{c44e0c8ee63ca166450922a373c7409c5d26b00b}. Each contains a
dependency manifest, exact primary-source snapshots, a readable proof,
source and prior-work checks, an independent audit, and two fresh build
records. The builds reuse dependency caches while rebuilding each
package's own proof objects from its delivered source.

\begin{center}
\begin{tabular}{@{}l>{\raggedright\arraybackslash}p{4.7in}@{}}
\toprule
Package & Main checked endpoints\\
\midrule
A326925 & \texttt{odd\_degree\_irreducible};
\texttt{even\_derivative\_numerator\_irreducible}; exact derivative and degree bridges.\\
A288668 & \texttt{word\_three\_steps}; \texttt{original\_conjecture}; literal rewriting, weighted recurrence, prefixes and growth.\\
A134559 & \texttt{witness\_fiber}; \texttt{witness\_prime\_pair}; \texttt{witness\_gcd}; \texttt{conjecture\_false}.\\
A000033 & \texttt{menage\_eq\_sourceMenage}; uniform polynomial shift; the three entry-specific divisibility theorems.\\
\bottomrule
\end{tabular}
\end{center}

The printed dependency reports use only Lean's standard axioms:
\texttt{propext}, \texttt{Classical.choice}, and \texttt{Quot.sound}, as
applicable. The proofs contain no admissions, extra axioms, unsafe
declarations, native evaluation shortcuts, or unproved experimental
recurrence assumptions. Finite computations in the counterexample use
ordinary proof-producing arithmetic or kernel reduction.

The source assertions and classical ingredients retain their original
attribution. The odd binomial criterion is an existing theorem, and the
m\'enage defining formula is classical. In particular, the proof of the
three source congruences is complete, but their historical novelty is
\emph{not established}: Carlitz's 1952 and 1954 papers and Stevens's 1981
paper are directly relevant, and their full mathematical statements were
unavailable to the review~\cite{Carlitz1952,Carlitz1954,Stevens1981}.
The accessible Wyman--Moser paper~\cite{WymanMoser1958} does not reproduce
the needed polynomial congruence.

Both search engines, the public trureturing index, and relevant nearby
definitions were checked for the other three targets; no exact earlier
resolution was located. An unchanged OEIS conjecture label and a bounded
search cannot certify that no earlier proof exists. Independent review
here means review by a separate collaborating agent and an isolated
rebuild; it does not mean external journal peer review or editorial
acceptance by OEIS. No OEIS submission is represented as completed.

\begin{thebibliography}{99}
\bibitem{A326925} C. Kimberling and OEIS contributors,
A326925, derivative-numerator coefficient array, 2019.
\url{https://oeis.org/A326925}. Exact revision 13, October 31, 2019.

\bibitem{A288665} C. Kimberling and OEIS contributors,
A288665, binary replacement process.
\url{https://oeis.org/A288665}. Exact revision 21, July 27, 2024.

\bibitem{A288668} C. Kimberling and OEIS contributors,
A288668, conjectured lengths of replacement iterates.
\url{https://oeis.org/A288668}. Exact revision 6, April 7, 2020.

\bibitem{A134559} G. W. Adamson, L. Ratajczak, and OEIS contributors,
A134559. Ratajczak's antidiagonal-product conjecture is dated
November 1, 2019. \url{https://oeis.org/A134559}.
Exact revision 16, September 20, 2025.

\bibitem{A005383} The OEIS Foundation Inc., A005383,
primes $p$ such that $(p+1)/2$ is prime.
\url{https://oeis.org/A005383}. Exact revision 143, August 4, 2026.

\bibitem{A000033} The OEIS Foundation Inc., A000033.
Prime-shift conjecture of M. van Hoeij, June 10, 2019.
\url{https://oeis.org/A000033}. Exact revision 60, June 21, 2026.

\bibitem{A000159} The OEIS Foundation Inc., A000159.
Prime-shift conjecture of M. van Hoeij, June 10, 2019.
\url{https://oeis.org/A000159}. Exact revision 54, June 21, 2026.

\bibitem{A000181} The OEIS Foundation Inc., A000181.
Prime-shift conjecture of M. van Hoeij, June 10, 2019.
\url{https://oeis.org/A000181}. Exact revision 30, June 21, 2026.

\bibitem{A058087} The OEIS Foundation Inc., A058087,
coefficients of the m\'enage hit polynomials.
\url{https://oeis.org/A058087}. Exact revision 66, November 5, 2025.
The defining formula is attributed to Kaplansky and Riordan.

\bibitem{Carlitz1952} L. Carlitz, Congruences for the m\'enage polynomials,
\emph{Duke Mathematical Journal} \textbf{19} (1952), 549--552.
\href{https://doi.org/10.1215/S0012-7094-52-01957-1}{doi:10.1215/S0012-7094-52-01957-1}.
Full text unavailable during this review.

\bibitem{Carlitz1954} L. Carlitz, Congruence properties of the m\'enage
polynomials, \emph{Scripta Mathematica} \textbf{20} (1954), 51--57.
Bibliographic details confirmed in~\cite{WymanMoser1958}; full text
unavailable during this review.

\bibitem{Stevens1981} H. Stevens, Arithmetic properties of the m\'enage
polynomials, \emph{Rocky Mountain Journal of Mathematics} \textbf{11}
(1981), 635--639.
\href{https://doi.org/10.1216/RMJ-1981-11-4-635}{doi:10.1216/RMJ-1981-11-4-635}.
Full text unavailable during this review; page range follows the JSTOR
issue contents.

\bibitem{WymanMoser1958} M. Wyman and L. Moser, On the probl\`eme des
m\'enages, \emph{Canadian Journal of Mathematics} \textbf{10} (1958),
468--480.
\href{https://doi.org/10.4153/CJM-1958-045-6}{doi:10.4153/CJM-1958-045-6}.

\bibitem{mathlib} The mathlib contributors. Odd binomial irreducibility,
theorem \texttt{X\_pow\_sub\_C\_irreducible\_of\_odd} in the
\texttt{Polynomial} namespace. Source file:
\path{Mathlib/FieldTheory/KummerPolynomial.lean}. Pinned revision:\par
\texttt{c44e0c8ee63ca166450922a373c7409c5d26b00b}.\par
\url{https://github.com/leanprover-community/mathlib4}.
\end{thebibliography}
\end{document}

